\FigureLayoutDeclare{img/2021/national_paper_a_science/q23_fig1.png}{12.0cm}{0bp 14bp 13bp 0bp}

\chapter{2021年全国甲卷（理）}

\section{单选题}

\begin{problem}
设集合 \(M=\{x\mid 0<x<4\}\)，\(N=\{x\mid \frac{1}{3}\le x\le 5\}\)，则 \(M\cap N=\)（\quad）

\choices
    {\(\left\{x\mid 0<x\le \frac{1}{3}\right\}\)}
    {\(\left\{x\mid \frac{1}{3}\le x<4\right\}\)}
    {\(\{x\mid 4\le x<5\}\)}
    {\(\{x\mid 0<x\le 5\}\)}

\answer{B}
\end{problem}

\begin{problem}
为了解某地农村经济情况，对该地农户家庭年收入进行抽样调查，将该地农户家庭年收入的调查数据整理得到如下频率分布直方图：

\bitmapfigure[width=0.75\linewidth]{img/2021/national_paper_a_science/q2_fig1.png}

根据此频率分布直方图，下面结论中不正确的是（\quad）

\choices
    {该地农户家庭年收入低于 \(4.5\text{ 万元}\) 的农户比率估计为 \(6\%\)}
    {该地农户家庭年收入不低于 \(10.5\text{ 万元}\) 的农户比率估计为 \(10\%\)}
    {估计该地农户家庭年收入的平均值不超过 \(6.5\) 万元}
    {估计该地有一半以上的农户，其家庭年收入介于 \(4.5\text{ 万元}\) 至 \(8.5\text{ 万元}\) 之间}

\answer{C}
\end{problem}

\begin{problem}
已知 \(\left(1-\i\right)^2z=3+2\i\)，则 \(z=\)（\quad）

\choices
    {\(-1-\frac{3}{2}\i\)}
    {\(-1+\frac{3}{2}\i\)}
    {\(-\frac{3}{2}+\i\)}
    {\(-\frac{3}{2}-\i\)}

\answer{B}
\end{problem}

\begin{problem}
青少年视力是社会普遍关注的问题，视力情况可借助视力表测量. 通常用五分记录法和小数记录法记录视力数据，五分记录法的数据 \(L\) 和小数记录法的数据 \(V\) 满足 \(L=5+\lg V\). 已知某同学视力的五分记录法的数据为 \(4.9\)，则其视力的小数记录法的数据约为（\(\sqrt[10]{10}\approx 1.259\)）（\quad）

\choices
    {\(1.5\)}
    {\(1.2\)}
    {\(0.8\)}
    {\(0.6\)}

\answer{C}
\end{problem}

\begin{problem}
已知 \(F_1\)，\(F_2\) 是双曲线 \(C\) 的两个焦点，\(P\) 为 \(C\) 上一点，且 \(\angle F_1PF_2=60^\circ\)，\(PF_1=3PF_2\)，则 \(C\) 的离心率为（\quad）

\choices
    {\(\dfrac{\sqrt{7}}{2}\)}
    {\(\dfrac{\sqrt{13}}{2}\)}
    {\(\sqrt{7}\)}
    {\(\sqrt{13}\)}

\answer{A}
\end{problem}

\begin{problem}
在一个正方体中，过顶点 \(A\) 的三条棱的中点分别为 \(E\)，\(F\)，\(G\). 该正方体截去三棱锥 \(A-EFG\) 后，所得多面体的三视图中，正视图如图所示，则相应的侧视图是（\quad）

\bitmapfigure[width=0.42\linewidth]{img/2021/national_paper_a_science/q6_fig1.png}

\choices
    {\choicebitmap[width=0.15\paperwidth]{img/2021/national_paper_a_science/q6_optA.png}}
    {\choicebitmap[width=0.15\paperwidth]{img/2021/national_paper_a_science/q6_optB.png}}
    {\choicebitmap[width=0.15\paperwidth]{img/2021/national_paper_a_science/q6_optC.png}}
    {\choicebitmap[width=0.15\paperwidth]{img/2021/national_paper_a_science/q6_optD.png}}

\answer{D}
\end{problem}

\begin{problem}
等比数列 \(\{a_n\}\) 的公比为 \(q\)，前 \(n\) 项和为 \(S_n\)，设甲：\(q>0\)，乙：\(\{S_n\}\) 是递增数列，则（\quad）

\choices
    {甲是乙的充分条件但不是必要条件}
    {甲是乙的必要条件但不是充分条件}
    {甲是乙的充要条件}
    {甲既不是乙的充分条件也不是乙的必要条件}

\answer{B}
\end{problem}

\begin{problem}
2020年12月8日，中国和尼泊尔联合公布珠穆朗玛峰最新高程为 \(8848.86\)（单位：\(\mathrm{m}\)），三角高程测量法是珠峰高程测量方法之一. 如图是三角高程测量法的一个示意图，现有 \(A\)，\(B\)，\(C\) 三点，且 \(A\)，\(B\)，\(C\) 在同一水平面上的投影 \(A'\)，\(B'\)，\(C'\) 满足 \(\angle A'C'B'=45^\circ\)，\(\angle A'B'C'=60^\circ\). 由 \(C\) 点测得 \(B\) 点的仰角为 \(15^\circ\)，\(BB'\) 与 \(CC'\) 的差为 \(100\)；由 \(B\) 点测得 \(A\) 点的仰角为 \(45^\circ\)，则 \(A\)，\(C\) 两点到水平面 \(A'B'C'\) 的高度差 \(AA'-CC'\) 约为（\(\sqrt{3}\approx 1.732\)）（\quad）

\bitmapfigure[width=0.45\linewidth]{img/2021/national_paper_a_science/q8_fig1.png}

\choices
    {\(346\)}
    {\(373\)}
    {\(446\)}
    {\(473\)}

\answer{B}
\end{problem}

\begin{problem}
若 \(\alpha\in \left( 0,\frac{\pi}{2} \right)\)，\(\tan 2\alpha=\dfrac{\cos \alpha}{2-\sin \alpha}\)，则 \(\tan \alpha=\)（\quad）

\choices
    {\(\dfrac{\sqrt{15}}{15}\)}
    {\(\dfrac{\sqrt{5}}{5}\)}
    {\(\dfrac{\sqrt{5}}{3}\)}
    {\(\dfrac{\sqrt{15}}{3}\)}

\answer{A}
\end{problem}

\begin{problem}
将 \(4\) 个 \(1\) 和 \(2\) 个 \(0\) 随机排成一行，则 \(2\) 个 \(0\) 不相邻的概率为（\quad）

\choices
    {\(\frac{1}{3}\)}
    {\(\frac{2}{5}\)}
    {\(\frac{2}{3}\)}
    {\(\frac{4}{5}\)}

\answer{C}
\end{problem}

\begin{problem}
已知 \(A\)，\(B\)，\(C\) 是半径为 \(1\) 的球 \(O\) 的球面上的三个点，且 \(AC\perp BC\)，\(AC=BC=1\)，则三棱锥 \(O-ABC\) 的体积为（\quad）

\choices
    {\(\dfrac{\sqrt{2}}{12}\)}
    {\(\dfrac{\sqrt{3}}{12}\)}
    {\(\dfrac{\sqrt{2}}{4}\)}
    {\(\dfrac{\sqrt{3}}{4}\)}

\answer{A}
\end{problem}

\begin{problem}
设函数 \(f(x)\) 的定义域为 \(\R\)，\(f(x+1)\) 为奇函数，\(f(x+2)\) 为偶函数，当 \(x\in [1,2]\) 时，\(f(x)=ax^2+b\). 若 \(f(0)+f(3)=6\)，则 \(f\left( \frac{9}{2} \right)=\)（\quad）

\choices
    {\(-\dfrac{9}{4}\)}
    {\(-\dfrac{3}{2}\)}
    {\(\dfrac{7}{4}\)}
    {\(\dfrac{5}{2}\)}

\answer{D}
\end{problem}

\section{填空题}

\begin{problem}
曲线 \(y=\dfrac{2x-1}{x+2}\) 在点 \((-1,-3)\) 处的切线方程为\fillinblank{}.

\answer{\(5x-y+2=0\)}
\end{problem}

\begin{problem}
已知向量 \(\bs a=(3,1)\)，\(\bs b=(1,0)\)，\(\bs c=\bs a+k\bs b\). 若 \(\bs a\perp \bs c\)，则 \(k=\fillinblank{}\).

\answer{\(-\dfrac{10}{3}\)}
\end{problem}

\begin{problem}
已知 \(F_1\)，\(F_2\) 为椭圆 \(C:\dfrac{x^2}{16}+\dfrac{y^2}{4}=1\) 的两个焦点，\(P\)，\(Q\) 为 \(C\) 上关于坐标原点对称的两点，且 \(PQ=F_1F_2\)，则四边形 \(PF_1QF_2\) 的面积为\fillinblank{}.

\answer{8}
\end{problem}

\begin{problem}
已知函数 \(f(x)=2\cos (\omega x+\varphi)\) 的部分图像如图所示，则满足条件 \(\left( f(x)-f\left( -\frac{7\pi}{4} \right) \right)\allowbreak\left( f(x)-f\left( \frac{4\pi}{3} \right) \right)>0\) 的最小正整数 \(x\) 为\fillinblank{}.

\bitmapfigure[width=0.5\linewidth]{img/2021/national_paper_a_science/q16_fig1.png}

\answer{2}
\end{problem}

\section{解答题}

\begin{problem}
甲，乙两台机床生产同种产品，产品质量分为一级品和二级品，为了比较两台机床产品的质量，分别用两台机床各生产了 \(200\) 件产品，产品的质量情况统计如下表：

\begin{center}
\begin{tabular}{c|ccc}
\hline
& 一级品 & 二级品 & 合计 \\
\hline
甲机床 & 150 & 50 & 200 \\
乙机床 & 120 & 80 & 200 \\
合计 & 270 & 130 & 400 \\
\hline
\end{tabular}
\end{center}

\begin{enumerate}
\item 甲机床，乙机床生产的产品中一级品的频率分别是多少？
\item 能否有 \(99\%\) 的把握认为甲机床的产品质量与乙机床的产品质量有差异？
\end{enumerate}

附：\(K^2=\dfrac{n(ad-bc)^2}{(a+b)(c+d)(a+c)(b+d)}\).

\begin{center}
\begin{tabular}{c|ccc}
\hline
\(P(K^2\ge k)\) & 0.050 & 0.010 & 0.001 \\
\hline
\(k\) & 3.841 & 6.635 & 10.828 \\
\hline
\end{tabular}
\end{center}
\end{problem}

\solution{
\begin{enumerate}
\item 甲机床生产的产品中一级品的频率为
\(\dfrac{150}{200}=0.75\)，乙机床生产的产品中一级品的频率为
\(\dfrac{120}{200}=0.60\).
\item 设原假设为 \(H_0\)：甲、乙两台机床生产的产品质量相互独立，即两台机床生产一级品的频率没有差异. 在附表公式中取
\(a=150\)，\(b=50\)，\(c=120\)，\(d=80\)，\(n=400\)，则
\[
K^2=\frac{400\left(150\times80-50\times120\right)^2}{200\times200\times270\times130}=\frac{400}{39}\approx10.256
\]
因为 \(10.256>6.635\)，所以在显著性水平 \(0.01\) 下拒绝原假设，即有 \(99\%\) 的把握认为甲机床的产品质量与乙机床的产品质量有差异.
\end{enumerate}
}

\begin{problem}
已知数列 \(\{a_n\}\) 的各项均为正数，记 \(S_n\) 为 \(\{a_n\}\) 的前 \(n\) 项和，从下面三个条件中选取两个作为条件，证明另外一个成立.

\begin{circlelist}
\item 数列 \(\{a_n\}\) 是等差数列；
\item 数列 \(\left\{\sqrt{S_n}\right\}\) 是等差数列；
\item \(a_2=3a_1\).
\end{circlelist}
\end{problem}

\solution{
\begin{enumerate}
\item 选择条件（1）和条件（3）. 设数列 \(\{a_n\}\) 的公差为 \(d\)，则
\(a_2=a_1+d=3a_1\)，从而 \(d=2a_1\). 由于各项均为正数，\(a_1>0\)，所以
\[
S_n=na_1+\frac{n(n-1)}2d=na_1+n(n-1)a_1=n^2a_1
\]
因此 \(\sqrt{S_n}=n\sqrt{a_1}\)，数列 \(\{\sqrt{S_n}\}\) 是等差数列.
\item 选择条件（1）和条件（2）. 设 \(\{a_n\}\) 的公差为 \(d\)，则
\[
S_n=na_1+\frac{n(n-1)}2d=\frac d2n^2+\left(a_1-\frac d2\right)n
\]
设 \(\{\sqrt{S_n}\}\) 的首项为 \(u\)，公差为 \(v\)，则 \(\sqrt{S_n}=u+(n-1)v\). 两边平方后，由于等式对所有正整数 \(n\) 成立，得到的两个关于 \(n\) 的多项式恒等. 左式在 \(n=0\) 时为 \(0\)，故右式的常数项满足 \((u-v)^2=0\)，即 \(u=v\). 于是 \(S_n=n^2v^2\). 比较 \(n^2\) 项和 \(n\) 项的系数，得 \(d=2v^2\)，\(a_1-d/2=0\)，所以 \(a_1=v^2\)，从而
\(a_2=a_1+d=v^2+2v^2=3a_1\)，条件（3）成立.
\item 选择条件（2）和条件（3）. 设 \(\{\sqrt{S_n}\}\) 的公差为 \(d\). 因为
\(S_1=a_1\)，\(S_2=a_1+a_2=4a_1\)，所以
\[
d=\sqrt{S_2}-\sqrt{S_1}=2\sqrt{a_1}-\sqrt{a_1}=\sqrt{a_1}
\]
即 \(a_1=d^2\). 因而 \(\sqrt{S_n}=\sqrt{S_1}+(n-1)d=nd\)，从而 \(S_n=n^2d^2\). 于是
\[
a_n=S_n-S_{n-1}=\left[n^2-(n-1)^2\right]d^2=(2n-1)d^2
\]
所以 \(\{a_n\}\) 是等差数列.
\end{enumerate}
}

\begin{problem}
已知直三棱柱 \(ABC-A_1B_1C_1\) 中，侧面 \(AA_1B_1B\) 为正方形，\(AB=BC=2\)，\(E\)，\(F\) 分别为 \(AC\) 和 \(CC_1\) 的中点，\(D\) 为棱 \(A_1B_1\) 上的点，\(BF\perp A_1B_1\).

\begin{enumerate}
\item 证明：\(BF\perp DE\)；
\item 当 \(B_1D\) 为何值时，面 \(BB_1C_1C\) 与面 \(DEF\) 所成的二面角的正弦值最小？
\end{enumerate}

\bitmapfigure[width=0.45\linewidth]{img/2021/national_paper_a_science/q19_fig1.png}
\end{problem}

\solution{
\begin{enumerate}
\item 由 \(A_1B_1\parallel AB\) 以及 \(BF\perp A_1B_1\)，并结合 \(F\) 是 \(CC_1\) 的中点，得 \(BF\perp AB\). 又 \(\overrightarrow{BF}=\overrightarrow{BC}+\dfrac12\overrightarrow{CC_1}\)，且 \(\overrightarrow{AB}\perp\overrightarrow{CC_1}\)，所以 \(\overrightarrow{AB}\cdot\overrightarrow{BC}=0\)，即 \(AB\perp BC\). 取 \(B\) 为原点，\(BA\)，\(BC\)，\(BB_1\) 分别为 \(x\)，\(y\)，\(z\) 轴正向，建立空间直角坐标系，则
\(A(2,0,0)\)，\(B(0,0,0)\)，\(C(0,2,0)\)，\(A_1(2,0,2)\)，\(B_1(0,0,2)\)，\(C_1(0,2,2)\). 因为 \(E\)，\(F\) 分别为 \(AC\)，\(CC_1\) 的中点，所以
\(E(1,1,0)\)，\(F(0,2,1)\). 设 \(B_1D=m\)，则 \(0\le m\le2\)，且 \(D(m,0,2)\). 因而
\(\overrightarrow{BF}=(0,2,1)\)，\(\overrightarrow{DE}=(1-m,1,-2)\).
所以 \(\overrightarrow{BF}\cdot\overrightarrow{DE}=2-2=0\)，即 \(BF\perp DE\).
\item 记面 \(BB_1C_1C\) 与面 \(DEF\) 所成的锐二面角为 \(\theta\). 面 \(BB_1C_1C\) 的一个法向量为 \(\bs{n}_1=(1,0,0)\). 又
\(\overrightarrow{DE}=(1-m,1,-2)\)，\(\overrightarrow{EF}=(-1,1,1)\)，
所以面 \(DEF\) 的一个法向量为
\(\bs{n}_2=\overrightarrow{DE}\times\overrightarrow{EF}=(3,m+1,2-m)\). 因此
\[
\sin^2\theta=1-\frac{\left(\bs{n}_1\cdot\bs{n}_2\right)^2}{|\bs{n}_1|^2|\bs{n}_2|^2}
 =\frac{(m+1)^2+(2-m)^2}{9+(m+1)^2+(2-m)^2}
 =\frac{4\left(m-\frac12\right)^2+9}{4\left(m-\frac12\right)^2+27}
\]
令 \(u=\left(m-\frac12\right)^2\ge0\)，则 \(\sin^2\theta=\dfrac{4u+9}{4u+27}\)，且该式随 \(u\) 增大而增大. 所以当 \(u=0\)，即 \(B_1D=m=\dfrac12\) 时，二面角的正弦值最小.
\end{enumerate}
}

\begin{problem}
抛物线 \(C\) 的顶点为坐标原点 \(O\)，焦点在 \(x\) 轴上，直线 \(l:x=1\) 交 \(C\) 于 \(P\)，\(Q\) 两点，且 \(OP\perp OQ\). 已知点 \(M(2,0)\)，且 \(\odot M\) 与 \(l\) 相切.

\begin{enumerate}
\item 求 \(C\)，\(\odot M\) 的方程；
\item 设 \(A_1\)，\(A_2\)，\(A_3\) 是 \(C\) 上的三个点，直线 \(A_1A_2\)，\(A_1A_3\) 均与 \(\odot M\) 相切. 判断直线 \(A_2A_3\) 与 \(\odot M\) 的位置关系，并说明理由.
\end{enumerate}
\end{problem}

\solution{
\begin{enumerate}
\item 设抛物线 \(C\) 的方程为 \(y^2=2px\)，其中 \(p>0\)，因为直线 \(x=1\) 与 \(C\) 有两个交点. 直线 \(x=1\) 与其交于 \(P\)，\(Q\)，故可设 \(P(1,\sqrt{2p})\)，\(Q(1,-\sqrt{2p})\). 由 \(OP\perp OQ\)，得
\(\overrightarrow{OP}\cdot\overrightarrow{OQ}=1-2p=0\)，所以 \(p=\dfrac12\)，即 \(C:y^2=x\). 圆 \(\odot M\) 的圆心为 \(M(2,0)\)，且与直线 \(x=1\) 相切，故半径等于圆心到直线的距离 \(1\)，从而 \(\odot M:(x-2)^2+y^2=1\).
\item 将抛物线 \(C\) 上的点参数化为 \(T(t)=(t^2,t)\). 设 \(A_1=T(t_1)\)，\(A_2=T(t_2)\)，\(A_3=T(t_3)\). 对抛物线上参数为 \(u\)，\(v\) 的两个不同点，连接这两点的直线方程为
\[
x-(u+v)y+uv=0
\]
圆心 \(M(2,0)\) 到该直线的距离为 \(1\) 当且仅当
\[
\frac{|2+uv|}{\sqrt{1+(u+v)^2}}=1
\]
即
\[
(u^2-1)v^2+2uv+3-u^2=0
\]
由于直线 \(A_1A_2\)，\(A_1A_3\) 均与圆相切，\(t_2\)，\(t_3\) 是关于 \(v\) 的方程
\((t_1^2-1)v^2+2t_1v+3-t_1^2=0\) 的两个根. 先说明 \(t_1^2\ne1\)：若 \(t_1=1\)，上式化为 \(2v+2=0\)；若 \(t_1=-1\)，上式化为 \(-2v+2=0\)，两种情形都不可能得到两个不同的根. 因此由韦达定理，记 \(s=t_2+t_3\)，\(r=t_2t_3\)，则
\(s=-\frac{2t_1}{t_1^2-1}\)，\(r=\frac{3-t_1^2}{t_1^2-1}\).
直线 \(A_2A_3\) 与圆相切的充要条件为 \((2+r)^2=1+s^2\)，等价于 \(r^2-s^2+4r+3=0\). 代入上式，得
\[
r^2-s^2+4r+3
=\frac{(3-t_1^2)^2-4t_1^2+4(3-t_1^2)(t_1^2-1)+3(t_1^2-1)^2}{(t_1^2-1)^2}
=\frac{0}{(t_1^2-1)^2}=0
\]
所以直线 \(A_2A_3\) 与 \(\odot M\) 相切.
\end{enumerate}
}

\begin{problem}
已知 \(a>0\) 且 \(a\ne 1\)，函数 \(f(x)=\dfrac{x^a}{a^x}\)（\(x>0\)）.

\begin{enumerate}
\item 当 \(a=2\) 时，求 \(f(x)\) 的单调区间；
\item 若曲线 \(y=f(x)\) 与直线 \(y=1\) 有且仅有两个交点，求 \(a\) 的取值范围.
\end{enumerate}
\end{problem}

\solution{
\begin{enumerate}
\item 当 \(a=2\) 时，\(f(x)=\dfrac{x^2}{2^x}\). 对 \(x>0\) 求导，得
\[
f'(x)=\frac{x(2-x\ln2)}{2^x}
\]
因为 \(2^x>0\) 且 \(x>0\)，所以当 \(0<x<\dfrac{2}{\ln2}\) 时 \(f'(x)>0\)，当 \(x>\dfrac{2}{\ln2}\) 时 \(f'(x)<0\). 故 \(f(x)\) 在 \(\left(0,\dfrac{2}{\ln2}\right)\) 上单调递增，在 \(\left(\dfrac{2}{\ln2},+\infty\right)\) 上单调递减.
\item 曲线与直线 \(y=1\) 的交点个数等于方程 \(f(x)=1\) 在 \(x>0\) 上的解的个数. 由 \(a>0\)，可取对数，得
\[
f(x)=1\Longleftrightarrow a\ln x=x\ln a\Longleftrightarrow \frac{\ln x}{x}=\frac{\ln a}{a}
\]
令 \(h(x)=\dfrac{\ln x}{x}\)，则 \(h'(x)=\dfrac{1-\ln x}{x^2}\). 因而 \(h(x)\) 在 \((0,\e)\) 上单调递增，在 \((\e,+\infty)\) 上单调递减，且最大值为 \(h(\e)=\dfrac1\e\). 当 \(0<a<1\) 时，\(\dfrac{\ln a}{a}<0\)，方程只有一个解；当 \(a>1\) 时，方程有两个解当且仅当
\(0<\dfrac{\ln a}{a}<\dfrac1\e\). 由于 \(\dfrac{\ln a}{a}\) 在 \((1,\e)\) 上递增，在 \((\e,+\infty)\) 上递减，故上述条件等价于 \(a\in(1,\e)\cup(\e,+\infty)\).
\end{enumerate}
}

\section{选考题：解答题}

\begin{problem}
在直角坐标系 \(xOy\) 中，以坐标原点为极点，\(x\) 轴正半轴为极轴建立极坐标系，曲线 \(C\) 的极坐标方程为 \(\rho=2\sqrt{2}\cos \theta\).

\begin{enumerate}
\item 将 \(C\) 的极坐标方程化为直角坐标方程；
\item 设点 \(A\) 的直角坐标为 \((1,0)\)，\(M\) 为 \(C\) 上的动点，点 \(P\) 满足 \(\overrightarrow{AP}=\sqrt{2}\overrightarrow{AM}\)，写出 \(P\) 的轨迹 \(C_1\) 的参数方程，并判断 \(C\) 与 \(C_1\) 是否有公共点.
\end{enumerate}
\end{problem}

\solution{
\begin{enumerate}
\item 由 \(\rho^2=x^2+y^2\)，\(x=\rho\cos\theta\)，原方程可得
\(\rho^2=2\sqrt2\rho\cos\theta=2\sqrt2x\). 因此曲线 \(C\) 的直角坐标方程为
\[
x^2+y^2=2\sqrt2x
\]
即 \(C:(x-\sqrt2)^2+y^2=2\).
\item 设 \(M\) 的参数坐标为 \(\left(\sqrt2+\sqrt2\cos\theta,\sqrt2\sin\theta\right)\)，其中 \(0\le\theta<2\pi\). 设 \(P(x,y)\)，由 \(A(1,0)\) 及
\(\overrightarrow{AP}=\sqrt2\overrightarrow{AM}\)，得
\[
\begin{cases}
x-1=\sqrt2\left(\sqrt2+\sqrt2\cos\theta-1\right),\\
y=\sqrt2\cdot\sqrt2\sin\theta.
\end{cases}
\]
所以 \(P\) 的轨迹 \(C_1\) 的参数方程为
\[
\begin{cases}
x=3-\sqrt2+2\cos\theta,\\
y=2\sin\theta.
\end{cases}
\]
故 \(C_1\) 是圆心为 \(O_1(3-\sqrt2,0)\)、半径为 \(2\) 的圆，而 \(C\) 是圆心为 \(O(\sqrt2,0)\)、半径为 \(\sqrt2\) 的圆. 两圆心距离为
\(OO_1=3-2\sqrt2<2-\sqrt2\)，即小圆半径与大圆半径之差大于两圆心距离，所以两圆内含且没有公共点.
\end{enumerate}
}

\begin{problem}
已知函数 \(f(x)=|x-2|\)，\(g(x)=|2x+3|-|2x-1|\).

\begin{enumerate}
\item 画出 \(y=f(x)\) 和 \(y=g(x)\) 的图像；
\item 若 \(f(x+a)\ge g(x)\)，求 \(a\) 的取值范围.
\end{enumerate}

\bitmapfigure[float=inline,max width=0.5\linewidth]{img/2021/national_paper_a_science/q23_fig1.png}
\FloatBarrier
\end{problem}

\solution{
\begin{enumerate}
\item
\(f(x)=\begin{cases}
2-x,&x<2,\\
x-2,&x\geqslant2,
\end{cases}\)
，\(g(x)=\begin{cases}
-4,&x<-\dfrac{3}{2},\\
4x+2,&-\dfrac{3}{2}\leqslant x<\dfrac{1}{2},\\
4,&x\geqslant\dfrac{1}{2},
\end{cases}\).
由此可知，\(y=f(x)\) 的图像是顶点为 \((2,0)\) 的折线，\(y=g(x)\) 的图像由左侧射线 \(y=-4\)、线段 \(y=4x+2\) 和右侧射线 \(y=4\) 组成，连接点分别为 \(\left(-\dfrac32,-4\right)\) 和 \(\left(\dfrac12,4\right)\).
\item \(a\) 取何值使不等式成立，指的是对所有 \(x\in\R\) 都有 \(f(x+a)\ge g(x)\). 当 \(x\ge\dfrac12\) 时，\(g(x)=4\)，所以必须有
\(\lvert x+a-2\rvert\ge4\) 对一切 \(x\ge\dfrac12\) 成立. 若 \(a\le\dfrac32\)，则点 \(x=2-a\) 属于 \(\left[\dfrac12,+\infty\right)\)，此时 \(f(x+a)=0<4\)，不合题意；若 \(a>\dfrac32\)，则 \(x+a-2\) 在该区间上递增且最小值为 \(a-\dfrac32\)，故必须 \(a-\dfrac32\ge4\)，即 \(a\ge\dfrac{11}{2}\).
反之，若 \(a\ge\dfrac{11}{2}\)，则当 \(x< -\dfrac32\) 时，\(g(x)=-4\le0\le f(x+a)\). 当 \(-\dfrac32\le x<\dfrac12\) 时，\(x+a-2\ge a-\dfrac72\ge2>0\)，从而
\[
f(x+a)-(4x+2)=x+a-2-(4x+2)=a-3x-4\ge\frac{11}{2}-3\cdot\frac12-4=0
\]
当 \(x\ge\dfrac12\) 时，\(x+a-2\ge a-\dfrac32\ge4\)，所以 \(f(x+a)\ge4=g(x)\). 因此不等式对所有 \(x\) 成立，且 \(a\) 的取值范围为 \(\left[\dfrac{11}{2},+\infty\right)\).
\end{enumerate}
}
